%!TEX TS-program = xelatex
%!TEX encoding = UTF-8 Unicode

% Awesome-CV vuelve a solicitar geometry al final de la clase con estas opciones.
% Pasarlas antes de cargar la clase evita el option clash; más abajo se aplican
% los márgenes definitivos mediante \geometry{...}.
\PassOptionsToPackage{left=5cm,top=2cm,right=1.5cm,bottom=2.5cm,nohead,nofoot}{geometry}
\documentclass[10pt,a4paper]{awesome-cv}
\addbibresource{bibliography.bib}

\fontdir[fonts/]
\colorlet{awesome}{awesome-red}
\setbool{acvSectionColorHighlight}{true}
\renewcommand{\acvHeaderSocialSep}{\quad\textbar\quad}

%-------------------------------------------------------------------------------
% INFORMACIÓN PERSONAL
%-------------------------------------------------------------------------------
\photo{profile.png}
\name{Jeremías}{Likerman}
\position{Investigador adjunto - CONICET{\enskip\cdotp\enskip} Profesor adjunto - Universidad de Buenos Aires}
\address{Instituto de Estudios Andinos ``Don Pablo Groeber'' (IDEAN), CONICET--UBA}

\mobile{(+54) 230-455-7758}
\email{jlikerman@uba.ar}
\googlescholar{HGA3sQQAAAAJ}{Google Scholar}
\extrainfo{\href{https://orcid.org/0000-0003-3844-2085}{ORCID 0000-0003-3844-2085}}

\usepackage[spanish]{babel}
% Shared Awesome-CV layout for the Spanish and English full CVs.
\usepackage{needspace}
\usepackage{tabularx}
\usepackage{longtable}
\usepackage{enumitem}
\usepackage{multicol}

\raggedbottom
\setlength{\emergencystretch}{2em}

\makeatletter

\renewcommand*{\sectionstyle}[1]{{\fontsize{14.2pt}{16.5pt}\bodyfont\bfseries\color{text}\@sectioncolor #1}}
\renewcommand*{\subsectionstyle}[1]{{\fontsize{10pt}{11.8pt}\bodyfont\bfseries\scshape\textcolor{graytext}{#1}}}
\renewcommand*{\paragraphstyle}{\fontsize{8.8pt}{11.3pt}\bodyfontlight\upshape\color{text}}
\renewcommand*{\entrytitlestyle}[1]{{\fontsize{9.5pt}{11.2pt}\bodyfont\bfseries\color{darktext} #1}}
\renewcommand*{\entrypositionstyle}[1]{{\fontsize{8pt}{9.5pt}\bodyfont\scshape\color{graytext} #1}}
\renewcommand*{\entrydatestyle}[1]{{\fontsize{8pt}{9.5pt}\bodyfontlight\slshape\color{graytext} #1}}
\renewcommand*{\entrylocationstyle}[1]{{\fontsize{8.3pt}{9.8pt}\bodyfontlight\slshape\color{awesome} #1}}
\renewcommand*{\descriptionstyle}[1]{{\fontsize{8.65pt}{11pt}\bodyfontlight\upshape\color{text} #1}}

\renewcommand{\cvsection}[1]{%
  \par\Needspace{4\baselineskip}%
  \addvspace{3.2mm}%
  \sectionstyle{#1}\par\nobreak%
  \vspace{0.15mm}%
  {\color{gray}\hrule height 0.55pt}%
  \vspace{0.65mm}%
}

\renewcommand{\cvsubsection}[1]{%
  \par\Needspace{3.5\baselineskip}%
  \addvspace{2mm}%
  \subsectionstyle{#1}\par\nobreak%
  \vspace{0.55mm}%
}

\renewenvironment{cvparagraph}{%
  \paragraphstyle
}{%
  \par\vspace{1mm}%
}

\renewenvironment{cventries}{%
  \par\vspace{0.1mm}%
}{%
  \par\vspace{0mm}%
}

\renewcommand*{\cventry}[5]{%
  \par\Needspace{3.8\baselineskip}%
  \addvspace{0mm}%
  \begingroup%
  \renewcommand{\arraystretch}{0.82}%
  \noindent\begin{tabularx}{\textwidth}{@{}X>{\raggedleft\arraybackslash}p{3.75cm}@{}}
    \entrytitlestyle{#2} & \entrylocationstyle{#3} \\
    \entrypositionstyle{#1} & \entrydatestyle{#4}
  \end{tabularx}%
  \endgroup%
  \vspace{-1.15mm}%
  \descriptionstyle{#5}\par%
  \vspace{-0.95mm}%
}

\renewenvironment{cvitems}{%
  \begin{itemize}[leftmargin=1.55em,topsep=0pt,itemsep=0pt,parsep=0pt,partopsep=0pt]
  \fontsize{8.65pt}{10.9pt}\bodyfontlight\color{text}
}{%
  \end{itemize}%
  \vspace{-1.1mm}%
}

\setlength{\LTpre}{0.4mm}
\setlength{\LTpost}{0.4mm}
\renewenvironment{cvhonors}{%
  \begin{longtable}{@{}p{1.35cm}p{\dimexpr\textwidth-4.25cm\relax}>{\raggedleft\arraybackslash}p{2.9cm}@{}}
}{%
  \end{longtable}%
}

\renewcommand*{\cvhonor}[4]{%
  \honordatestyle{#4} &
  \raggedright\honorpositionstyle{#1}\par\honortitlestyle{#2} &
  \honorlocationstyle{#3}\tabularnewline[0.9mm]
}

\makeatother

\DeclareFieldFormat[article]{title}{#1}
\defbibenvironment{bibliography}
  {\begin{enumerate}[leftmargin=1.6em,label=\arabic*.,font=\fontsize{7.55pt}{8.45pt}\bodyfont\mdseries\color{darktext},topsep=0pt,itemsep=0.15mm,parsep=0pt,partopsep=0pt]}
  {\end{enumerate}}
  {\item}
\DeclareBibliographyDriver{article}{%
  {\fontsize{7.55pt}{8.45pt}\bodyfont\mdseries\color{darktext}\printfield{title}}%
  \setunit{\addperiod\space}%
  {\fontsize{6.95pt}{7.8pt}\selectfont\color{graytext}\printnames{author}}%
  \setunit{\addperiod\space}%
  {\fontsize{6.9pt}{7.75pt}\selectfont\addfontfeature{Color=lightgray}\itshape%
    \usebibmacro{journal+issuetitle}%
    \setunit{\space}%
    \printfield{pages}%
  }%
  \par\vspace{0.12mm}%
}


\begin{document}

\makecvheader

\makecvfooter
  {\today}
  {Jeremías Likerman~~~·~~~Curriculum Vitae}
  {\thepage}

\vspace{1.2mm}
%-------------------------------------------------------------------------------
%	SECTION TITLE
%-------------------------------------------------------------------------------
\cvsection{Perfil}

\begin{cvparagraph}

\textbf{Perfil de investigación:} Geólogo especializado en geología estructural, geomecánica y modelado numérico multiescala, con más de 15 años de experiencia en investigación, docencia universitaria, trabajo de campo y formación de estudiantes. Su investigación integra modelos bidimensionales y tridimensionales para estudiar procesos geodinámicos, geotérmicos y geomecánicos en los Andes y en reservorios energéticos, incluyendo zonas de subducción, deformación cortical, sistemas volcánicos con potencial geotérmico y fracturación natural e inducida, con énfasis en la Formación Vaca Muerta.

\textbf{Docencia, transferencia y gestión:} Profesor adjunto de la Universidad de Buenos Aires e investigador adjunto del CONICET. Su trayectoria docente universitaria comenzó en 2013 como Ayudante de Primera, continuó como Jefe de Trabajos Prácticos entre 2018 y 2026 y prosigue como Profesor Adjunto desde 2026. Ha dictado Geología Estructural, Geotectónica, Tectónica de Campo, Introducción a la Geología, Levantamiento Geológico, Geoestadística y Geoinformática; dirige y codirige investigadores y estudiantes; coordina proyectos nacionales e internacionales y participa en actividades de transferencia, divulgación científica y gestión institucional en el IDEAN.

\end{cvparagraph}

%-------------------------------------------------------------------------------
%	SECTION TITLE
%-------------------------------------------------------------------------------
\cvsection{Formación académica y posdoctoral}


%-------------------------------------------------------------------------------
%	CONTENT
%-------------------------------------------------------------------------------
\begin{cventries}

%---------------------------------------------------------
%  \cventry
%    {B.S. in Computer Science and Engineering} % Degree
%    {POSTECH(Pohang University of Science and Technology)} % Institution
%    {Pohang, S.Korea} % Location
%    {Mar. 2010 - Aug. 2017} % Date(s)
%    {
%      \begin{cvitems} % Description(s) bullet points
%        \item {Got a Chun Shin-Il Scholarship which is given to promising students in CSE Dept.}
%      \end{cvitems}
%    }



\cventry
    {CONICET y UPC}
    {Beca posdoctoral}
    {Argentina/España}
    {2015 - 2018}
    {
      \begin{cvitems} % Description(s) bullet points
        \item {Tema de investigación: evolución térmica de la litosfera oceánica.
Directores: Dr. Cristallini (CONICET-UBA) y Dr. Zlotnik (UPC).}
      \end{cvitems}
    }

\cventry
    {Stanford University}
    {Investigador visitante en Stanford}
    {Estados Unidos}
    {2013}
    {
      \begin{cvitems} % Description(s) bullet points
        \item {\emph{Cursos realizados: Geología Estructural y MATLAB aplicado a Geociencias.}}
      \end{cvitems}
    }

\cventry
    {Universidad de Buenos Aires}
    {Doctorado en Ciencias Geológicas}
    {Argentina}
    {2010 - 2015}
    {
      \begin{cvitems} % Description(s) bullet points
        \item {Tesis: Predicción del fracturamiento en superficies geológicas plegadas mediante modelado numérico. Director: Dr. Cristallini.}
      \end{cvitems}
    }

\cventry
    {Universidad de Buenos Aires}
    {Licenciatura en Ciencias Geológicas}
    {Argentina}
    {2005 - 2010}
    {
      \begin{cvitems} % Description(s) bullet points
        \item {Tesis: Estructura y geología del anticlinal compuesto Tres Cruces, Jujuy, Argentina. Director: Dr. Cristallini.}
      \end{cvitems}
    }

%---------------------------------------------------------
\end{cventries}

%-------------------------------------------------------------------------------
% SECTION TITLE
%-------------------------------------------------------------------------------
\cvsection{Cargos y experiencia profesional}

\cvsubsection{Cargos actuales}

\begin{cventries}

\cventry
    {CONICET--IDEAN}
    {Investigador adjunto}
    {Argentina}
    {2018--presente}
    {
      \begin{cvitems}
        \item {Carrera del Investigador Científico desde 2018; promoción a Investigador Adjunto en noviembre de 2023. Líneas de trabajo: predicción de fracturas, modelado geodinámico y evaluación numérica del potencial geotérmico.}
      \end{cvitems}
    }

\cventry
    {Departamento de Ciencias Geológicas, Universidad de Buenos Aires}
    {Profesor adjunto}
    {Argentina}
    {Junio de 2026--presente}
    {
      \begin{cvitems}
        \item {Asignaturas: Geoestadística, Geoinformática y Levantamiento Geológico. Dedicación exclusiva.}
      \end{cvitems}
    }

\cventry
    {IDEAN, CONICET--UBA}
    {Coordinador institucional}
    {Argentina}
    {Septiembre de 2025--presente}
    {
      \begin{cvitems}
        \item {Integrante del Consejo Directivo del instituto desde abril de 2024.}
      \end{cvitems}
    }

\end{cventries}

\Needspace{9\baselineskip}
\cvsubsection{Cargos docentes previos}

\begin{cventries}

\cventry
    {Departamento de Ciencias Geológicas, Universidad de Buenos Aires}
    {Jefe de Trabajos Prácticos}
    {Argentina}
    {Marzo de 2018--mayo de 2026}
    {
      \begin{cvitems}
        \item {Asignaturas: Introducción a la Geología, Levantamiento Geológico y Geoestadística.}
      \end{cvitems}
    }

\cventry
    {Departamento de Ciencias Geológicas, Universidad de Buenos Aires}
    {Ayudante de Primera}
    {Argentina}
    {Junio de 2013--septiembre de 2017}
    {
      \begin{cvitems}
        \item {Asignaturas: Geología Estructural, Geotectónica y Tectónica de Campo.}
      \end{cvitems}
    }

\end{cventries}

\cvsubsection{Cursos y actividades de campo}

\begin{cventries}

\cventry
    {Docente a cargo}
    {Geología de reservorios no convencionales}
    {Argentina}
    {2025}
    {
      \begin{cvitems}
        \item {Curso de formación profesional orientado a la caracterización geológica de reservorios de shale.}
      \end{cvitems}
    }

\cventry
    {Organizador e instructor}
    {Salida avanzada de investigación sobre Vaca Muerta}
    {Argentina}
    {2025}
    {
      \begin{cvitems}
        \item {\emph{Perspectivas sobre Vaca Muerta: atributos del subsuelo analizados desde superficie}.}
      \end{cvitems}
    }

\end{cventries}

\cvsubsection{Estadías de investigación}

\begin{cventries}

\cventry
    {GEO3BCN--CSIC}
    {Investigador visitante}
    {España}
    {Mayo--septiembre de 2025}
    {
      \begin{cvitems}
        \item {Modelado termomecánico de zonas volcánicas para la evaluación del potencial geotérmico. Anfitriona: Dra. Ivone Jiménez-Munt.}
      \end{cvitems}
    }

\cventry
    {GEO3BCN--CSIC}
    {Investigador visitante}
    {España}
    {Agosto--septiembre de 2024}
    {
      \begin{cvitems}
        \item {Desarrollo de modelos numéricos acoplados para sistemas volcánicos y geotérmicos.}
      \end{cvitems}
    }

\cventry
    {BarcelonaTech (UPC)}
    {Investigador visitante}
    {España}
    {2015, 2017, 2019 y 2022}
    {
      \begin{cvitems}
        \item {Estadías de investigación en el Laboratorio de Cálculo Numérico.}
      \end{cvitems}
    }

\end{cventries}

%-------------------------------------------------------------------------------
% SECTION TITLE
%-------------------------------------------------------------------------------
\cvsection{Proyectos de investigación}

\begingroup
\renewcommand*{\honorpositionstyle}[1]{{\fontsize{9pt}{1em}\bodyfont\mdseries\color{darktext} #1}}
\begin{cvhonors}

\cvhonor
    {Análisis tectónico-estructural mediante modelado análogo y numérico: contribuciones a los recursos energéticos (co-investigador responsable).}
    {PIP 11220250100817CO, CONICET}
    {Argentina}
    {2026--2029}

\cvhonor
    {Modelos tridimensionales de transferencia de esfuerzos entre segmentos de flat slab andinos (investigador responsable).}
    {AECT-2026-1-0080}
    {España}
    {2026}

\cvhonor
    {Potencial geotérmico de zonas volcánicas (investigador responsable).}
    {AECT-2025-1-0009}
    {España}
    {2025}

\cvhonor
    {Deformación cuaternaria en el Sistema de Famatina (co-investigador).}
    {Fundación Williams}
    {Argentina}
    {2025--2026}

\cvhonor
    {Estimación del potencial geotérmico en zonas volcánicas mediante modelos numéricos termomecánicos (investigador responsable).}
    {LINCGL2024}
    {España}
    {2024--2026}

\cvhonor
    {Modelado de la geodinámica de zonas de subducción y plumas mantélicas (investigador responsable).}
    {AECT-2024-2-0027}
    {España}
    {2024}

\cvhonor
    {Traverse structures in the Andean orogen (integrante).}
    {PICT 2021-I-A-01210}
    {Argentina}
    {2023--2027}

\cvhonor
    {Geodinámica de la microplaca Adria, desde el manto hasta la superficie (co-investigador responsable).}
    {GEO3BCN}
    {España}
    {2023}

\cvhonor
    {Modelado geodinámico de zonas de subducción (investigador responsable).}
    {AECT-2022-3-0023}
    {España}
    {2022}

\cvhonor
    {Estudio de la dinámica del paisaje y de los riesgos geológicos asociados al emplazamiento de presas en áreas montañosas (co-investigador responsable).}
    {PICT-2019-04011}
    {Argentina}
    {2020}

\cvhonor
    {Modelado análogo y numérico: contribuciones a la resolución de problemas tectónicos y estructurales (co-investigador responsable).}
    {PICT-2019-00997}
    {Argentina}
    {2020}

\cvhonor
    {Fracturación natural y estructuras de primer orden: modelos predictivos aplicados a Cerro Domuyo (co-investigador responsable).}
    {UBACyT}
    {Argentina}
    {2018}

\cvhonor
    {Evolución geológica de los Andes e impacto económico y ambiental (co-investigador responsable).}
    {PUE 22920160100051}
    {Argentina}
    {2017}

\cvhonor
    {Modelado numérico de la evolución térmica de la litosfera oceánica (investigador responsable).}
    {PICT 2015-1226}
    {Argentina}
    {2016}

\cvhonor
    {Aplicación de modelos análogos y numéricos a problemas tectónico-estructurales (co-investigador responsable).}
    {PICT-2016-1407}
    {Argentina}
    {2016}

\end{cvhonors}
\endgroup

%-------------------------------------------------------------------------------
%	SECTION TITLE
%-------------------------------------------------------------------------------
\newpage
\cvsection{Dirección de recursos humanos}

\cvsubsection{Dirección de investigadores}

\begin{cventries}

%---------------------------------------------------------
\cventry
    {Investigadora asistente} % Role
    {Plotek, Berenice} % Event
    {Argentina} % Location
    {Desde agosto de 2025} % Date(s)
    {
      \vspace{-0.8mm}
      \begin{cvitems} % Description(s)
        \item {\textbf{Director:} Likerman, J. \textbf{Codirector:} Cristallini, E.}
      \end{cvitems}
    }
\end{cventries}

\vspace{3.5mm}
\cvsubsection{Dirección posdoctoral}

\begin{cventries}

%---------------------------------------------------------
\cventry
    {Pliegues por propagación de falla: aproximación mecánica mediante modelado numérico por elementos finitos} % Role
    {Plotek, Berenice} % Event
    {Argentina} % Location
    {2024} % Date(s)
    {
      \vspace{-0.8mm}
      \begin{cvitems} % Description(s)
        \item {\textbf{Dirección:} Cristallini, E. y Likerman, J.}
      \end{cvitems}
    }
\end{cventries}

\vspace{3.5mm}
\cvsubsection{Dirección de tesis doctorales}

%-------------------------------------------------------------------------------
%	CONTENT
%-------------------------------------------------------------------------------
\begin{cventries}

%---------------------------------------------------------
  \cventry
    {Fracturación natural en la Formación Vaca Muerta: caracterización y modelado} % Role
    {Correa Luna, Clara} % Event
    {Argentina} % Location
    {2023} % Date(s)
    {
      \vspace{-0.8mm}
      \begin{cvitems} % Description(s)
        \item {\textbf{Dirección:} Yagupsky, D. y Likerman, J.}
      \end{cvitems}
    }
\end{cventries}

\vspace{3.5mm}
\cvsubsection{Dirección de tesis de licenciatura}
\begin{cventries}

% \cventry
%     {Quaternary deformation in the Chaschuil Valley, Western slope of the Sierra de la Planchadas, Catamarca}
%     {Gramajo, Pilar}
%     {Argentina} % Location
%     {In progress} % Date(s)
%     {
%       \begin{cvitems} % Description(s)
%         \item {\textbf{Advisors:} Sagripanti, L. y Likerman, J.}
%       \end{cvitems}
%     }

\cventry
    {Inversión tectónica de un complejo metamórfico en el Macizo del Deseado: aportes desde el modelado análogo}
    {Chandía Sepúlveda, Alexander Ignacio}
    {Argentina} % Location
    {En curso} % Date(s)
    {
      \vspace{-0.8mm}
      \begin{cvitems} % Description(s)
        \item {\textbf{Dirección:} Likerman, J. y Navarrete, C.}
      \end{cvitems}
    }

\cventry
    {Deformación cuaternaria en el valle de Fiambalá, Catamarca}
    {De Rosa, Tomás}
    {Argentina} % Location
    {En curso} % Date(s)
    {
      \vspace{-0.8mm}
      \begin{cvitems} % Description(s)
        \item {\textbf{Dirección:} Likerman, J. y Sagripanti, L.}
      \end{cvitems}
    }


\cventry
    {Análisis del potencial geotermoeléctrico del Complejo Volcánico Cerro Blanco, Catamarca: influencia de la reología y estructura a través de modelos numéricos.}
    {Blanco, Lucía}
    {Argentina} % Location
    {En curso} % Date(s)
    {
      \vspace{-0.8mm}
      \begin{cvitems} % Description(s)
        \item {\textbf{Dirección:} Likerman, J. y Lamberti, C.}
      \end{cvitems}
    }


\cventry
    {Evolución estructural de un complejo metamórfico invertido en el Macizo del Deseado: aportes desde el modelado numérico}
    {Celis Santisteban, Celso Angel}
    {Argentina} % Location
    {En curso} % Date(s)
    {
      \vspace{-0.8mm}
      \begin{cvitems} % Description(s)
        \item {\textbf{Dirección:} Likerman, J. y Navarrete, C.}
      \end{cvitems}
    }

\cventry
    {Evaluación del potencial geotermoeléctrico del Complejo Volcánico Cerro Blanco (Catamarca) mediante modelado numérico.}
    {Ruiz, Luciana}
    {Argentina} % Location
    {2026} % Date(s)
    {
      \vspace{-0.8mm}
      \begin{cvitems} % Description(s)
        \item {\textbf{Dirección:} Lamberti, C. y Likerman, J.}
      \end{cvitems}
    }

\cventry
    {Deformación cuaternaria en el valle de Chaschuil, vertiente occidental de la Sierra de Las Planchadas, provincia de Catamarca}
    {Gramajo, Pilar}
    {Argentina} % Location
    {2025} % Date(s)
    {
      \vspace{-0.8mm}
      \begin{cvitems} % Description(s)
        \item {\textbf{Dirección:} Likerman, J. y Sagripanti, L.}
      \end{cvitems}
    }



%---------------------------------------------------------
  \cventry
    {Geomorfología del valle del río Grande entre Bardas Blancas y Portezuelo del Viento, Mendoza}
    {Acosta, Luana}
    {Argentina} % Location
    {2023} % Date(s)
    {
      \vspace{-0.8mm}
      \begin{cvitems} % Description(s)
        \item {\textbf{Dirección:} Winocur, D. y Likerman, J.}
      \end{cvitems}
    }

%---------------------------------------------------------
  \cventry
     {Geología de la ladera oriental del Cerro Perito Moreno, El Bolsón, provincia de Río Negro}
    {Mazzini, Martín} % Event
    {Argentina} % Location
    {2022} % Date(s)
    {
      \vspace{-0.8mm}
      \begin{cvitems} % Description(s)
        \item {\textbf{Dirección:} Likerman, J. y Tobal, J.}
      \end{cvitems}
    }

%---------------------------------------------------------

  \cventry
     {Estructura y neotectónica del área de Tabladitas, provincia de Jujuy, Argentina}
    {Relañez, Gastón} % Event
    {Argentina} % Location
    {2019} % Date(s)
    {
      \vspace{-0.8mm}
      \begin{cvitems} % Description(s)
        \item {\textbf{Dirección:} Yagupsky, D. y Likerman, J.}
      \end{cvitems}
    }

%---------------------------------------------------------

  \cventry
     {Geología, estructura y fracturación de la Formación Lajas en el anticlinal Covunco, Portada Covunco, Neuquén}
    {Ventisky, Esteban} % Event
    {Argentina} % Location
    {2017} % Date(s)
    {
      \vspace{-0.8mm}
      \begin{cvitems} % Description(s)
        \item {\textbf{Dirección:} Likerman, J. y Tobal, J.}
      \end{cvitems}
    }

%---------------------------------------------------------

  \cventry
     {Geología estructural y fracturación en la cuenca de Tres Cruces, provincia de Jujuy}
    {Correa Luna, Clara} % Event
    {Argentina} % Location
    {2015} % Date(s)
    {
      \vspace{-0.8mm}
      \begin{cvitems} % Description(s)
        \item {\textbf{Dirección:} Yagupsky, D. y Likerman, J.}
      \end{cvitems}
    }

%---------------------------------------------------------

\end{cventries}

%-------------------------------------------------------------------------------
%	SECTION TITLE
%-------------------------------------------------------------------------------
\cvsection{Becas y distinciones}


%-------------------------------------------------------------------------------
%	SUBSECTION TITLE
%-------------------------------------------------------------------------------
%\cvsubsection{International}


%-------------------------------------------------------------------------------
%	CONTENT
%-------------------------------------------------------------------------------
\begin{cvhonors}

%---------------------------------------------------------

\cvhonor
    {Beca de investigación en Japón} % Role
    {Kobe University} % Event
    {JSPS} % Location
    {2027} % Date(s)

\cvhonor
    {Viaje de campo SZNet 2025 Chile} % Role
    {National Science Foundation} % Event
    {SZ4D} % Location
    {2025} % Date(s)

\cvhonor
    {Becas de Movilidad Académica entre Instituciones Asociadas} % Role
    {Asociación Universitaria Iberoamericana de Posgrado} % Event
    {AUIP} % Location
    {2022} % Date(s)
    

\cvhonor
    {Consejo Nacional de Investigaciones Científicas y Técnicas}
    {Estadía internacional de investigación}
    {CONICET}
    {2016}
    

\cvhonor
    {Consejo Nacional de Investigaciones Científicas y Técnicas}
    {Beca Posdoctoral}
    {CONICET}
    {2015}
    

%---------------------------------------------------------
\cvhonor
    {Stanford University}
    {Fulbright - Fundación Bunge \& Born}
    {Beca doctoral}
    {2013}
    

%---------------------------------------------------------
\cvhonor
   {Consejo Nacional de Investigaciones Científicas y Técnicas}
    {Beca Doctoral Tipo II}
    {CONICET}
    {2013}
    

%---------------------------------------------------------
\cvhonor
  {Consejo Nacional de Investigaciones Científicas y Técnicas}
    {Beca Doctoral Tipo I}
    {CONICET}
    {2010}
    


% \cvhonor
%     {3D thermomechanical evolution of oceanic lithosphere.}
%     {AECT-2016-3-0010}
%     {Spain} % Location
%     {2016} % Date(s)

% %---------------------------------------------------------

% \cvhonor
%     {Modeling the evolution of orogenic wedges. Influence of erosion on tectonic activity.}
%     {PICT 2012-02176}
%     {Argentina} % Location
%     {2014} % Date(s)

% %---------------------------------------------------------

% \cvhonor
%     {Analog and numerical tectonic-structural modeling.}
%     {PICT 2013-1309}% Event
%     {Argentina} % Location
%     {2013} % Date(s)

% %---------------------------------------------------------

% \cvhonor
%     {Analog and numerical modeling of fault-related folds.}
%     {UBACyT X855} % Award% Event
%     {Argentina} % Location
%     {2011} % Date(s)

% %---------------------------------------------------------

% \cvhonor
%     {Applicability of analog and numerical models to tectonic problems.}
%     {PICT-2010-1441} % Award% Event
%     {Argentina} % Location
%     {2010} % Date(s)

%---------------------------------------------------------

\end{cvhonors}


%-------------------------------------------------------------------------------
%	SUBSECTION TITLE
%-------------------------------------------------------------------------------
% \cvsubsection{Domestic}


% %-------------------------------------------------------------------------------
% %	CONTENT
% %-------------------------------------------------------------------------------
% \begin{cvhonors}

% %---------------------------------------------------------
%   \cvhonor
%     {3rd Place} % Award
%     {WITHCON Hacking Competition Final} % Event
%     {Seoul, S.Korea} % Location
%     {2015} % Date(s)

% %---------------------------------------------------------
%   \cvhonor
%     {Silver Prize} % Award
%     {KISA HDCON Hacking Competition Final} % Event
%     {Seoul, S.Korea} % Location
%     {2017} % Date(s)

% %---------------------------------------------------------
%   \cvhonor
%     {Silver Prize} % Award
%     {KISA HDCON Hacking Competition Final} % Event
%     {Seoul, S.Korea} % Location
%     {2013} % Date(s)

% %---------------------------------------------------------
%   \cvhonor
%     {2nd Award} % Award
%     {HUST Hacking Festival} % Event
%     {S.Korea} % Location
%     {2013} % Date(s)

% %---------------------------------------------------------
%   \cvhonor
%     {3rd Award} % Award
%     {HUST Hacking Festival} % Event
%     {S.Korea} % Location
%     {2010} % Date(s)

% %---------------------------------------------------------
%   \cvhonor
%     {3rd Award} % Award
%     {Holyshield 3rd Hacking Festival} % Event
%     {S.Korea} % Location
%     {2012} % Date(s)

% %---------------------------------------------------------
%   \cvhonor
%     {2nd Award} % Award
%     {Holyshield 3rd Hacking Festival} % Event
%     {S.Korea} % Location
%     {2011} % Date(s)

% %---------------------------------------------------------
%   \cvhonor
%     {5th Place} % Award
%     {PADOCON Hacking Competition Final} % Event
%     {Seoul, S.Korea} % Location
%     {2011} % Date(s)

% %---------------------------------------------------------
% \end{cvhonors}

%-------------------------------------------------------------------------------
% SECTION TITLE
%-------------------------------------------------------------------------------
\cvsection{Competencias}

\setlength{\columnsep}{6mm}
\begin{multicols}{2}
\raggedcolumns

\descriptionstyle{\textbf{Programación y reproducibilidad:} \LaTeX{}, Python, R, MATLAB, HTML5, Git y documentación científica reproducible.}\par
\vspace{0.3mm}

\descriptionstyle{\textbf{Modelado numérico y HPC:} Underworld2, ASPECT, GOLEM y LaMEM; MPI, posprocesamiento y visualización de modelos 2D/3D.}\par
\vspace{0.3mm}

\descriptionstyle{\textbf{Modelado análogo:} diseño experimental, caracterización de materiales, PIV y análisis cinemático y mecánico.}\par

\columnbreak

\descriptionstyle{\textbf{Análisis espacial y geológico:} QGIS, GMT y PostGIS; cartografía, interpretación sísmica, análisis estructural y métodos de campo asistidos por UAV.}\par
\vspace{0.3mm}

\descriptionstyle{\textbf{Trabajo de campo:} más de 15 años de experiencia en los Andes; planificación, logística y liderazgo de equipos.}\par
\vspace{0.3mm}

\descriptionstyle{\textbf{Idiomas:} español nativo, inglés fluido (C2), italiano (A2) y catalán (A2).}\par

\end{multicols}

%-------------------------------------------------------------------------------
% SECTION TITLE
%-------------------------------------------------------------------------------
\cvsection{Publicaciones con referato}

\sloppy
\begin{refsection}
  \nocite{*}
  \setlength{\bibitemsep}{0pt}
  \setlength{\bibparsep}{0pt}
  \printbibliography[sorting=chronological,type=article,heading=none]
\end{refsection}


\end{document}
